\chapter{Zakończenie}\label{ch:zakonczenie}

Celem pracy była kompleksowa analiza porównawcza samodzielnie zarządzanego klastra Valkey uruchomionego w~Kubernetes z~klastrem Redis~7.2 oraz z~usługą zarządzaną Google Cloud Memorystore for Redis, prowadzona pod kątem wydajności, opóźnień, odporności na awarie, kosztów oraz nakładu pracy operacyjnej. Postawiony problem badawczy --- czy Valkey może stanowić praktyczną, wydajną i~ekonomicznie uzasadnioną alternatywę dla Redis~7.2 i~rozwiązania zarządzanego --- doczekał się odpowiedzi opartej na pomiarach w~siedmiu kategoriach testów, powtarzanych pięciokrotnie i~poddanych analizie statystycznej z~korektą na wielokrotne porównania.

Zgromadzony materiał pozwala stwierdzić, że Valkey jest dojrzałą i~funkcjonalnie zgodną alternatywą dla Redis~7.2. W~większości profili obciążenia osiąga istotnie wyższą przepustowość, przy payloadach pośrednich (10~KB) różnice są statystycznie nieistotne, a~jedyny przypadek pozornej przewagi Redis (odczyt dużych wartości) okazał się artefaktem współczynnika trafień, a~nie cechą silnika. Pod względem zachowania podczas awarii mastera oraz spójności danych w~trakcie partycji sieciowej oba systemy są statystycznie nieodróżnialne, co jest spójne ze~wspólnym modelem asynchronicznej replikacji. Z~kolei atomowa migracja slotów dostępna w~nowszych wersjach Valkey istotnie skraca czas reshardingu i~ogranicza jego wpływ na klienta. Z~perspektywy ekonomicznej infrastruktura self-hosted okazała się tańsza od usługi zarządzanej, kosztem wyższego nakładu pracy operacyjnej.

Rezultaty pracy mają zarówno charakter poznawczy, jak i~użytkowy. W~warstwie poznawczej pokazano, że deklarowana zgodność interfejsu między Valkey a~Redis przekłada się na zgodność obserwowanego zachowania w~scenariuszach produkcyjnych, oraz że pozornie spektakularne różnice w~surowych metrykach mogą wynikać z~konfiguracji i~sposobu pomiaru, a~nie z~architektury silnika. W~warstwie użytkowej praca dostarcza zestawu rekomendacji oraz powtarzalnej metodyki i~narzędzi benchmarkowych, które organizacja stojąca przed decyzją migracyjną po zmianie licencji Redis może wykorzystać do oceny własnego obciążenia.

Przeprowadzone badania mają jednak ograniczenia, które wyznaczają granice wniosków. Badania wydajności prowadzono w~środowisku lokalnym (minikube) ze~względu na ograniczony dostęp do zasobów chmurowych, co zawyża wartości bezwzględne przepustowości i~utrudnia bezpośrednie przeniesienie ich do środowiska produkcyjnego. Liczba powtórzeń ($N=5$) jest wystarczająca do testów nieparametrycznych, ale ogranicza moc wykrywania drobnych różnic. Dedykowanego, pełnego testu wydajnościowego usługi Memorystore nie udało się zrealizować z~powodu ograniczeń zasobowych --- jej przepustowość oszacowano jedynie z~jednego profilu zarejestrowanego w~teście reshardingu, z~dodatkowym narzutem skoku sieciowego klienta. Porównanie samodzielnie zarządzanych klastrów i~usługi zarządzanej nie odbywa się też w~identycznym środowisku, a~różnice między użytymi chartami Helm (valkey-helm oraz Bitnami) wprowadzają asymetrie konfiguracji, których nie da się w~pełni wyeliminować.

Naturalnym kierunkiem dalszych prac jest powtórzenie testów wydajnościowych w~jednorodnym środowisku chmurowym o~większej skali, z~pełną macierzą profili obciążenia również dla Memorystore, oraz zwiększenie liczby powtórzeń w~celu uzyskania węższych przedziałów ufności. Wartościowe byłoby również zbadanie wpływu trwałości danych (AOF, snapshoty), szyfrowania TLS oraz dłuższych, bardziej złożonych scenariuszy chaos engineering na wnioski dotyczące niezawodności. Osobnym wątkiem jest domknięcie analizy zaobserwowanej asymetrii eksmisji przez porównanie efektywnych konfiguracji silników na działających instancjach, a~także automatyzacja operacji utrzymaniowych za pomocą operatorów Kubernetes, co mogłoby zmniejszyć przewagę usług zarządzanych w~obszarze nakładu pracy operacyjnej.
